\begin{problem}[Warmup]{(a)}
    Given a subgroup $H$ of a group $G$, at least one left coset of $H$ always equals its corresponding right coset. Why?
\end{problem}
\unsolved

\begin{problem}[Warmup]{(b)}
    Why is every subgroup normal, if the group is abelian? (The converse is false; see Exercise 10 below.)
\end{problem}
\unsolved

\begin{problem}[Warmup]{(c)}
    If $H$ is a normal subgroup of the group $G$, and $[G : H] = 7$, to what group is $G/H$ necessarily isomorphic?
\end{problem}
\unsolved

\begin{problem}[Warmup]{(d)}
    The ring $\mathbb{Z}$ is a subring of $\mathbb{R}$ but is not an ideal. (Why?) We thus were unable to form the ring of cosets $\mathbb{R}/\mathbb{Z}$. But $\mathbb{Z}$ is clearly a normal subgroup, and so we are able to form the group of cosets $\mathbb{R}/\mathbb{Z}$. Explain this. (For more information about the group $\mathbb{R}/\mathbb{Z}$, see Exercise 33.1.)
\end{problem}
\unsolved

\begin{problem}[Warmup]{(e)}
    Show that $\{1\}$ is a normal subgroup of any multiplicative group $G$.
\end{problem}
\unsolved

\begin{problem}[Warmup]{(f)}
    Consider the set \[\mathcal{K} = \{A \in U(M_2(\mathbb{R})) : \det(A) = 3\}.\] By using an argument just like that in Example 32.6, it is easy to see that if $A \in \mathcal{K}$, then $C^{-1}AC \in \mathcal{K}$, for all matrices $C \in U(M_2(\mathbb{R}))$. Does this make $\mathcal{K}$ a normal subgroup of $U(M_2(\mathbb{R}))$? (Be careful!)
\end{problem}
\unsolved

\begin{problem}{1}
    Suppose that $H$ is a normal subgroup of the group $G$. Prove that
    $H \subseteq g^{-1}Hg$. \emph{Note:} This is an omitted verification in the proof of
    Theorem 32.2.
\end{problem}
\unsolved

\begin{problem}{2}
    Suppose that $H$ is a subgroup of the group $G$, and $g^{-1}Hg = H$, for all $g \in G$.
    Prove that $Hg \subseteq gH$. \emph{Note:} This is another omitted verification in the
    proof of Theorem 32.2.
\end{problem}
\unsolved

\begin{problem}{3}
    Are $\langle (124) \rangle$ and $\langle (12) \rangle$ normal subgroups of $S_4$?
    (Consider Exercise 31.2.) Is $\langle (124) \rangle$ a normal subgroup of $A_4$?
    (Consider Exercise 31.1.)
\end{problem}
\unsolved

\begin{problem}{4}
    Is $\{\iota, (123), (132)\}$ a normal subgroup of $S_3$? (Consider Example 31.3.)
\end{problem}
\unsolved

\begin{problem}{5}
    Suppose that $n > 2$ is a positive integer and $1 \le k \le n$. Show that the
    stabilizer subgroup $G_k$ is not normal in $S_n$. (See Exercise 29.6.)
\end{problem}
\unsolved

\begin{problem}{6}
    Let $H$ be a (not necessarily normal) subgroup of the group $G$, and $g \in G$. Prove
    that $g^{-1}Hg$ is always a subgroup of $G$. What group is it, if $H$ is normal? The
    subgroup $g^{-1}Hg$ is called a \textbf{conjugate} of the subgroup $H$. Argue that if
    $H_1$ is a conjugate of the subgroup $H_2$, then $H_2$ is a conjugate of $H_1$ (that
    is, conjugacy is a \emph{symmetric} relation).
\end{problem}
\unsolved

\begin{problem}{7}
    Suppose that $n$ is a positive integer, and $d$ is a positive integer that (properly)
    divides $n$. Prove that
    \[\mathbb{Z}_n/\langle d \rangle\]
    is isomorphic to $\mathbb{Z}_d$. \emph{Note:} This exercise foreshadows developments in
    the next chapter and will be easy to do with the main theorem of that chapter; your
    experience with rings should suggest the appropriate approach. See Exercise 33.4.
\end{problem}
\unsolved

\begin{problem}{8}
    Suppose that $H$ and $K$ are normal subgroups of the group $G$. Prove that $H \cap K$
    is a normal subgroup. \emph{Note:} In Exercise 25.5 you have already shown that
    $H \cap K$ is a subgroup.
\end{problem}
\unsolved

\begin{problem}{9}
    Consider the direct product $G \times H$ of the groups $G$ and $H$. Prove that the
    subgroups $G \times \{1\}$ and $\{1\} \times H$ are normal subgroups of $G \times H$.
\end{problem}
\unsolved

\begin{problem}{10}
    Show that every subgroup of the group $Q_8$ of quaternions is normal, even though the
    group is not abelian. (See Warm-up Exercise b.)
\end{problem}
\unsolved

\begin{problem}{11}
    Consider the subgroup $H = \{1, -1\}$ of $Q_8$. In the previous exercise, you showed
    that $H$ is a normal subgroup of $Q_8$. Give the group table for $Q_8/H$. To what
    group is this group of cosets isomorphic?
\end{problem}
\unsolved

\begin{problem}{12}
    Consider the group of cosets $\mathbb{R}^*/\mathbb{Q}^*$. Exhibit an element of this
    group with order 2. What about an element of order $n$, for any positive integer $n$?
    Do you believe that there are elements of infinite order in this group? (You should be
    able to conjecture an element that works, but you will not be able to prove this
    rigorously.)
\end{problem}
\unsolved

\begin{problem}{13}
    Suppose that $H$ and $K$ are subgroups of the group $G$, and $K$ is a normal subgroup.
    Let $HK = \{hk : h \in H, k \in K\}$.
    \begin{enumerate}
        \item[(a)] Prove that $HK$ is a subgroup of $G$.
        \item[(b)] Suppose in addition that $H$ is normal. Prove that $HK$ is a normal
        subgroup of $G$.
    \end{enumerate}
\end{problem}
\unsolved

\begin{problem}{14}
    Suppose that $G$ is a group with subgroups $H, K$. Suppose further that $H$ is a normal
    subgroup and $H \subseteq K \subseteq G$. Prove that $H$ is a normal subgroup of the
    group $K$.
\end{problem}
\unsolved

\begin{problem}{15}
    Suppose $G$ is a group, and that $g, h \in G$. Furthermore, assume that the order of
    $h$ is $n$. Prove that the order of $g^{-1}hg$ is $n$. The element $g^{-1}hg$ is called
    a \textbf{group conjugate} of the element $h$; an element of group shares many
    properties with its conjugates.
\end{problem}
\unsolved

\begin{problem}{16}
    Consider the groups $A_4$ and $H$ described in Example 32.5. Use Exercise 15 to provide
    a complete proof that $H$ is normal in $A_4$.
\end{problem}
\unsolved

\begin{problem}{17}
    Let $G$ be an abelian group, and let
    \[t(G) = \{g \in G : o(g) \text{ is finite}\}.\]
    We call $t(G)$ the \textbf{torsion subgroup} of $G$.
    \begin{enumerate}
        \item[(a)] Prove that $t(G)$ is a subgroup of $G$.
        \item[(b)] Prove that $G/t(G)$ is torsion-free. (Recall that this means that $G$
        has no non-identity elements of finite order.)
    \end{enumerate}
\end{problem}
\unsolved

\begin{problem}{18}
    Suppose that $H$ is a normal subgroup of the group $G$. Prove that $G/H$ is an abelian
    group if and only if $g^{-1}h^{-1}gh \in H$, for all $g, h \in G$. \emph{Note:}
    Elements of the form $g^{-1}h^{-1}gh$ are called \textbf{commutators} in the group $G$.
\end{problem}
\unsolved

\begin{problem}{19}
    We generalize Exercise 18 slightly. Given a group $G$, define $H$ to be the smallest
    subgroup of $G$ that contains all the commutators. The subgroup $H$ is called the
    \textbf{commutator subgroup}. Argue that $H$ consists of all finite products of
    commutators. For a normal subgroup $K$ of $G$, prove that $G/K$ is abelian if and only
    if $K \supseteq H$.
\end{problem}
\unsolved
